\setcounter{chapter}{12} % Sets chapter counter so the current chapter is 13
\renewcommand{\thetheorem}{\thechapter.\arabic{theorem}}
\setcounter{theorem}{0}

\chapter{Row and Column Spaces}

In this chapter, we bridge the abstract theory of vector spaces with the algorithmic power of matrix theory. We use Gauss elimination as a scalpel to dissect the structure of coordinate spaces.

\subsection{Linear Combinations and Matrix Equations}

Consider the vector space $K^{m}$ over a field $K$. Let $w_{1}, \dots, w_{n} \in K^{m}$ be a collection of vectors. We begin by addressing three fundamental questions that govern our understanding of these spaces.

\begin{question}[The Three Fundamental Questions]
\label{q:foundational_questions}
\begin{enumerate}
    \item How can we systematically determine if a given vector $w \in K^{m}$ belongs to the span $\Sp(w_{1}, \dots, w_{n})$?
    \item Can we describe the subspace $\Sp(w_{1}, \dots, w_{n})$ precisely using algebraic conditions?
    \item How can we determine whether the vectors $w_{1}, \dots, w_{n}$ are linearly independent?
\end{enumerate}
\end{question}

\begin{answer}[To Questions 1 and 2]
By definition, $w \in \Sp(w_{1}, \dots, w_{n})$ if and only if there exist scalars $x_{1}, \dots, x_{n} \in K$ such that $w = x_{1} w_{1} + \dots + x_{n} w_{n}$. We can assemble the vectors $w_{j}$ as columns of an $m \times n$ matrix $A$. This abstract linear combination is equivalent to the matrix equation:
\[
\underbrace{\begin{pmatrix} | & & | \\ w_{1} & \dots & w_{n} \\ | & & | \end{pmatrix}}_{A} \cdot \begin{pmatrix} x_{1} \\ \vdots \\ x_{n} \end{pmatrix} = \begin{pmatrix} | \\ w \\ | \end{pmatrix}
\]
Thus, $w$ is in the span if and only if the linear system $Ax = w$ is \qt{consistent} (i.e., it possesses at least one valid solution vector $x$). We analyze this by applying Gauss elimination on the augmented matrix $(A \mid w)$. 

To answer the second question and describe the entire span precisely, we replace the specific target $w$ with a general, variable vector $b = (b_{1}, \dots, b_{m})^{T}$ and reduce the augmented matrix $(A \mid b)$ to a row-reduced echelon form $(A' \mid b')$. 

This row operation procedure transforms our original problem into a much simpler, equivalent linear system $A'x = b'$. As shown in Prof. Biran's notes, the resulting matrix generally has the following characteristic block structure:
\begin{equation}
\label{eq:block_consistency}
(A' \mid b') = \left( 
\begin{array}{ccc|c}
* & \dots & * & b'_{1} \\
\vdots & \text{\footnotesize Rows with pivots} & \vdots & \vdots \\
* & \dots & * & b'_{r} \\ \hline
0 & \dots & 0 & b'_{r+1} \\
\vdots & \text{\footnotesize Zero rows} & \vdots & \vdots \\
0 & \dots & 0 & b'_{m}
\end{array} 
\right)
\end{equation}
In this reduced block structure, any row of $A'$ consisting entirely of zeros translates to an algebraic equation of the form $0 = b'_{i}$. For the overall system to be consistent (meaning a solution actually exists), it is an absolute requirement that the right-hand side of these specific equations also evaluates to zero. 

In other words, the system is consistent if and only if the entries in $b'$ corresponding to the zero-rows of $A'$ strictly vanish. Therefore, we can concisely describe the span as the set of all vectors satisfying these constraints:
\[
\Sp(w_{1}, \dots, w_{n}) = \left\{ b \in K^{m} \;\middle|\; b'_{r+1} = \dots = b'_{m} = 0 \right\}.
\]
\end{answer}

% ex:13.1
\begin{example}[Calculating a Specific Span]
\label{ex:span_calc}
Let $v_{1} = (1, 1, 1)^{T}$ and $v_{2} = (1, 0, -1)^{T} \in K^{3}$. To describe $\Sp(v_{1}, v_{2})$, we set up the augmented matrix with a general vector $b$ and row-reduce:
\[
\left( \begin{array}{cc|c} 1 & 1 & b_{1} \\ 1 & 0 & b_{2} \\ 1 & -1 & b_{3} \end{array} \right) \xrightarrow{R_{2}-R_{1}, R_{3}-R_{1}} \left( \begin{array}{cc|c} 1 & 1 & b_{1} \\ 0 & -1 & b_{2}-b_{1} \\ 0 & -2 & b_{3}-b_{1} \end{array} \right) \xrightarrow{R_{3}-2R_{2}} \left( \begin{array}{cc|c} 1 & 1 & b_{1} \\ 0 & -1 & b_{2}-b_{1} \\ \hline 0 & 0 & b_{1}-2b_{2}+b_{3} \end{array} \right)
\]
The system is consistent if and only if the expression corresponding to the zero-row evaluates to zero. Thus, the span is perfectly described by that single plane equation:
\[
\Sp(v_{1}, v_{2}) = \{ (b_{1}, b_{2}, b_{3})^{T} \in K^{3} \mid b_{1} - 2b_{2} + b_{3} = 0 \}.
\]
\end{example}

\begin{answer}[To Question 3]
The vectors are linearly independent if and only if the homogeneous equation $Ax = 0$ has only the trivial solution (i.e., $x_1 = \dots = x_n = 0$). 

This occurs if and only if the row-reduced echelon form $A'$ has a pivot in every column. If even a single column lacks a pivot, the corresponding variable becomes \qt{free} (e.g., it can be parameterized to take on any scalar value). The existence of a free variable immediately allows for non-trivial solutions, which by definition means the vectors are linearly dependent.
\end{answer}

% thm:13.2
\begin{theorem}[Fundamental Constraints on Coordinate Vectors]
\label{thm:coordinate_constraints}
Let $v_{1}, \dots, v_{n} \in K^{m}$.
\begin{enumerate}
    \item If $n > m$, then the vectors $v_{1}, \dots, v_{n}$ are linearly dependent.
    \item If $n < m$, then the vectors $v_{1}, \dots, v_{n}$ cannot span $K^{m}$.
\end{enumerate}
\end{theorem}

\begin{proof}
\begin{enumerate}
    \item Let $A = (v_{1} \mid \dots \mid v_{n})$. The number of pivots in $A'$ is bounded by the number of rows $m$. If $n > m$, at least $n - m$ columns are forced to be without pivots. These free variables guarantee non-trivial solutions to $Ax = 0$, hence dependence.
    \item The number of pivots is bounded by the number of columns $n$. If $n < m$, there are at least $m - n$ zero-rows in $A'$. By the logic of equation \eqref{eq:block_consistency}, we can easily choose a vector $b$ such that the resulting $b'$ has non-zero entries in those bottom rows, making the system inconsistent.
\end{enumerate}
\end{proof}

\subsection{Row and Column Spaces}

% def:13.4
\begin{definition}[Subspaces of a Matrix]
\label{def:matrix_subspaces}
Let $A \in \M_{m \times n}(K)$.
\begin{enumerate}
\item The \qt{column space} $\ColS(A) \subseteq K^{m}$ is the span of the columns of $A$.
\item The \qt{row space} $\RowS(A) \subseteq K^{n}$ is the span of the rows of $A$.
\end{enumerate}
\end{definition}

% prop:13.5
\begin{proposition}[Invariance of Row Space]
\label{prop:row_space_invariance}
Elementary row operations do not change the row space: $\RowS(A) = \RowS(A')$.
\end{proposition}

\begin{proof}
New rows are, simply, linear combinations of the old rows, and because elementary operations are reversible, the spans remain identical.
\end{proof}

\begin{importantremark}
Row operations \textbf{do} change the column space, as seen by $\begin{pmatrix} 1 & 1 \\ 1 & 1 \end{pmatrix} \to \begin{pmatrix} 1 & 1 \\ 0 & 0 \end{pmatrix}$. However, they \textbf{preserve} the linear relations between columns: $\sum \alpha_{i} w_{i} = 0 \iff \sum \alpha_{i} w'_{i} = 0$.
\end{importantremark}

\subsection{Constructing Bases}

% thm:13.7
\begin{theorem}[Basis for the Column Space]
\label{thm:basis_col_space}
The columns of $A$ that correspond to the \qt{pivot columns} of its row-reduced echelon form $B$ constitute a basis for $\ColS(A)$.
\end{theorem}

\begin{proof}
By the principle that row operations preserve the linear dependence relations between columns, we have:
\[
\sum_{i=1}^n \alpha_i w_i = 0 \iff \sum_{i=1}^n \alpha_i w'_i = 0.
\]
Thus, a set of columns in $A$ is linearly independent (or spans the column space) if and only if the corresponding set of columns in $B$ possesses the same property in $\ColS(B)$.

Consider the row-reduced echelon matrix $B \in \M_{m \times n}(K)$ with $r$ pivots. We visualize its block structure and the specific indices of its pivot columns $j_1, j_2, \dots, j_r$:

\begin{equation}
\label{eq:matrix_B_rigorous}
B = \left( 
\begin{array}{ccccccccc}
 & \text{\tiny col } j_1 & & \text{\tiny col } j_2 & & \text{\tiny col } j_r & & & \\
 & \downarrow & & \downarrow & & \downarrow & & & \\
\dots & \mathbf{1} & * & \mathbf{0} & * & \mathbf{0} & * & \dots & \leftarrow 1 \\
\dots & \mathbf{0} & \dots & \mathbf{1} & \dots & \mathbf{0} & * & \dots & \leftarrow 2 \\
 & \vdots & & \vdots & \ddots & \vdots & \vdots & & \vdots \\
\dots & \mathbf{0} & \dots & \mathbf{0} & \dots & \mathbf{1} & * & \dots & \leftarrow r \\ \hline
 & 0 & & 0 & & 0 & 0 & \dots & \\
 & \vdots & & \vdots & & \vdots & \vdots & \ddots & \\
 & 0 & & 0 & & 0 & 0 & \dots & 
\end{array} 
\right)
\end{equation}

In this echelon form, the pivot columns are precisely the first $r$ standard basis vectors of $K^m$:
\[
e_1 = \begin{pmatrix} 1 \\ 0 \\ \vdots \\ 0 \\ \vdots \\ 0 \end{pmatrix}, \quad 
e_2 = \begin{pmatrix} 0 \\ 1 \\ \vdots \\ 0 \\ \vdots \\ 0 \end{pmatrix}, \quad \dots, \quad 
e_r = \begin{pmatrix} 0 \\ \vdots \\ 1 \\ 0 \\ \vdots \\ 0 \end{pmatrix}
\]
where for each $e_i$, the entry $1$ is located at the $i$-th row. Because $B$ is in row-reduced echelon form, every entry in every column $w'_j$ that lies below the $r$-th row must be zero. This implies that:
\[
w'_j \in K^r \times \{0\}^{m-r} = \left\{ (x_1, \dots, x_r, 0, \dots, 0)^{T} \in K^m \right\}.
\]
Since the standard basis vectors $e_1, \dots, e_r$ are clearly linearly independent and they span the subspace $K^r \times \{0\}^{m-r}$, we establish the following chain of inclusions:
\[
K^r \times \{0\}^{m-r} = \Sp(e_1, \dots, e_r) \subseteq \ColS(B) \subseteq K^r \times \{0\}^{m-r}.
\]
This double inclusion forces the equality $\ColS(B) = \Sp(e_1, \dots, e_r) = K^r \times \{0\}^{m-r}$. Consequently, the pivot columns of $B$ form a basis for its column space. By the preservation of linear relations, the corresponding columns in the original matrix $A$ form a basis for $\ColS(A)$.
\end{proof}

% thm:13.8
\begin{theorem}[Basis for the Row Space]
\label{thm:row_basis}
The non-zero rows of the row-reduced echelon matrix $B$ form a basis for $\RowS(A)$.
\end{theorem}

\begin{corollary}[Rank Invariance]
\label{cor:rank_invariance}
The number of pivots is independent of the sequence of row operations, as it always equals $\dim \RowS(A)$.
\end{corollary}

% thm:13.13
\begin{theorem}[The Rank Equality]
\label{thm:rank_equality}
For any $A \in \M_{m \times n}(K)$, $\dim \RowS(A) = \dim \ColS(A)$. This shared dimension is the \qt{rank} of $A$, denoted $\rank(A)$.
\end{theorem}

\begin{proof}
Let $r$ be the number of pivots in the RREF of $A$. By Theorem \ref{thm:basis_col_space}, $\dim \ColS(A) = r$. By Theorem \ref{thm:row_basis}, $\dim \RowS(A) = r$. Thus, the dimensions are exactly equal.
\end{proof}